\newpage
\section{Vermittlungsschicht}
  \subsection{Paketvermittlung}
    \subsubsection{Paketvermittlung (Packet Switching)}
      \begin{itemize}
        \item Paketvermittlung werden die Daten in Pakete geglieder und zu den Zielsystemen
          weitergeleitet. Hierzu führen die Pakete Metainformationen (Kontrollinformationenen) mit sich.
      \end{itemize} 
    \subsubsection{Eigenschaften}
      \begin{itemize}
        \item Abschnittsweise Übermittlung von Paketen (Hop-by-Hop)
        \item Zwischenspeicherung (Pufferung) erfolgt ggf. in Warteschlangen/Puffern
          (engl. Queue/Buffer) der Zwischensysteme (Führt zu variablen Verzögerungen)
        \item Verlust von Paketen möglich (Begrenzte Pufferkapazität in den Zwischensystemen)
      \end{itemize}
  \subsection{Prinzipien}
    \subsubsection{Vermittlungsschicht}
      \begin{itemize}
        \item Datenaustausch zwischen Endsystemen, die nicht direkt physiklaisch benachbart sind, d.h.
          die nicht mit dem gleichen physiklaischen Medium verbunden sind
          \begin{itemize}
            \item Pakete werden typischerweise als Datagramm bezeichnet
            \item Zwischensysteme werden als Router bezeichnet
          \end{itemize}
      \end{itemize}
    \subsubsection{Zwei grundlegende Aufgaben}
      \begin{itemize}
        \item Wegewahl (Routing) - Kontrollebene (Kontrollpfade)
          \begin{itemize}
            \item Ermittelt den Weg (Pfad), den die Pakete zurücklegen
            \item Erfordert Routingalgorithmus und -protokoll
          \end{itemize}
        \item Weiterleitung (Forwarding) - Datenebene (Datenpfad)
          \begin{itemize}
            \item Leitet Pakete vom Eingangsinterface des Routers an ein geeignetes Ausgangsinterface
              weiter
          \end{itemize}
      \end{itemize}
    \subsubsection{Routing und Weiterleitung}
      \begin{itemize}
        \item Wege für Paketweiterleitung zum Ziel finden
          \begin{itemize}
            \item In Netzen (Nutzung von Switches)
            \item Zwischen Netzten (Nutzung von Routern)
            \item Weiterleitungstabelle durch Routingprotokoll etabliert
          \end{itemize}
      \end{itemize}
    \subsubsection{Aufgabe der Kontrollebende: Wegewahl}
      \begin{itemize}
        \item Wegewahl (Routing) ermittelt den Weg (Pfad), den die Pakete vom Quell-Endsystem zum
          Ziel-Endsystem zurücklegen
        \item Weg (Pfad) setzt sich aus Sequenz von Routern zusammen. Diese sind über
          übertragungsabschnitte (Hops/Links) miteinander verbunden
          \begin{itemize}
            \item Übertragungsabschnitte
              \begin{itemize}
                \item Endsystem und benachbarter Router
                \item Router und benachbarter Router
                \item Router und benachbartes Endsystem
              \end{itemize}
          \end{itemize}
        \item Typischerweise
          \begin{itemize}
            \item Erster Router auf dem Pfad ist Default-Router für Quell-Endsystem
            \item Routingalgorithmen bestimmen Pfad zwischen Default-Router und Ziel-Router
            \item Ziel-Router leitet Paket direkt an Ziel-Endsystem weiter
          \end{itemize}
      \end{itemize}
    \subsubsection{Kontrollebene}
      \begin{itemize}
        \item Routingprotokoll
          \begin{itemize}
            \item Protokoll, das für die Wegewahl verantwortlich ist
            \item Oberhalb der Schicht 3 (Vermittlungsschicht) angesiedelt
          \end{itemize}
        \item Routingalgorithmus
          \begin{itemize}
            \item Algorithmus der für die Wegewahl eingesetzt wird
            \item Routingprotokolle setzen Routingalgorithmen in Rechnernetzten um
          \end{itemize}
        \item Weiterleitungstabelle
          \begin{itemize}
            \item Enthält Einträge für mögliche Ziele mit Interfaces, auf denen Datagramme in Richtung
              Ziel weitergeleitet werden
            \item Ist Schnittstelle zwischen Kontroll- und Datenebene
            \item Einträge in der Weiterleitungstabelle werden
              \begin{itemize}
                \item von Routingprotokollen geliefert
                \item in der Datenebene für die Weiterleitung der Datagramme ausgewertet
              \end{itemize}
          \end{itemize}
      \end{itemize}
    \subsubsection{Begriffe}
      \begin{itemize}
        \item Router
          \begin{itemize}
            \item Auf Vermittlungsschicht operierendes Zwischensystem
            \item Leitet eingehende Datagramme in Richtung Ziel weiter
              \begin{itemize}
                \item Nutzt Informationen in der Weiterleitungstabelle
              \end{itemize}
            \item Tauscht über Routingprotokolle Informationen mit anderen Routern aus
              \begin{itemize}
                \item Erstellung der Routing- und Weiterleitungstabellen bzw. Flow-Tables (SDN)
                  \begin{itemize}
                    \item Verteiltes adaptives Routing (traditionell) oder
                    \item logisch zentrales Routing (SDN)
                  \end{itemize}
              \end{itemize}
          \end{itemize}
        \item Route
          \begin{itemize}
            \item Hier: Weg (Pfad) eines Datagramms zum Zielsystem
              \begin{itemize}
                \item Daher auch "Wege"-wahl für die Auswahl der besten Route
              \end{itemize}
          \end{itemize}
        \item Link
          \begin{itemize}
            \item Hier: Übertragungsabschnitt zwischen zwei Routern
          \end{itemize}
        \item Port
          \begin{itemize}
            \item Eingangs- oder Ausgangsinterface
          \end{itemize}
      \end{itemize}
  \subsection{Datenebene}
    \subsubsection{IP-Adressen und Subnetze}
      IP-Adressen
      \begin{itemize}
        \item Ziel
          \begin{itemize}
            \item Eindeutige Identifikation aller Interfaces
              \begin{itemize}
                \item Ein Gerät kann mehrere Interfaces haben (z.B. Ethernet, WLAN, LTE)
              \end{itemize}
          \end{itemize}
        \item IP-Adressen
          \begin{itemize}
            \item Eindeutige Kennung für Interfaces von Routern und Endsystemen
            \item Einfaches, für Maschinen leicht zu verarbeitendes Format
            \item IPv4
              \begin{itemize}
                \item Adressen einer Länge von 32 Bit
              \end{itemize}
            \item IPv6
              \begin{itemize}
                \item Größerer Adressraum durch Adressen einer Länge von 128 Bit
              \end{itemize}
          \end{itemize}
      \end{itemize}
      IP-Adressen: Notation und Beispiel
      \begin{itemize}
        \item IPv4-Adresse (32 Bit)
          \begin{itemize}
            \item 4 Gruppen mit jeweils 8 Bit, separiert durch Punkt
          \end{itemize}
          \includegraphics[width=.8\textwidth]{img/ipv4.png}
        \item IPv6-Adresse (128 Bit)
          \begin{itemize}
            \item 8 Gruppen mit jeweils 16 Bit, separiert durch Doppelpunkt
              \begin{itemize}
                \item Jede Gruppe als 4 hexadezimale Zeichen notiert
              \end{itemize}
            \item Vereinfachungen
              \begin{itemize}
                \item Führende Nullen in jeder Gruppe werden entfernt
                \item Eine Riehe von Nullen kann durch :: ersetzt werden
                  \begin{itemize}
                    \item Typischerweise die längste Reihe von Nullen
                  \end{itemize}
              \end{itemize}
          \end{itemize}
          \includegraphics[width=.8\textwidth]{img/ipv6.png}
      \end{itemize}
      Gliederung von IPv4-Adressen
      \begin{itemize}
        \item Gliederung in
          \begin{itemize}
            \item Subnetz-Teil
              \begin{itemize}
                \item Kann unterschiedliche Länge haben, identifiziert das Netz
              \end{itemize}
            \item Endsystem-Teil
              \begin{itemize}
                \item Teil der Adresse die nicht zum Subnetz-Teil gehört, identifiziert Endsystem in
                  diesem Netz
              \end{itemize}
            \item $Format a.b.c.d/x$
              \begin{itemize}
                \item $x$: Anzahl der Bits im Subnetz-Teil
              \end{itemize}
          \end{itemize}
        \item Beispiel (IPv4) 200.63.16.0/23\\\\
          \includegraphics[width=.8\textwidth]{img/ipv4_bsp.png}
        \item Subnetzmaske
          \begin{itemize}
            \item "1"-Bits definieren Subnetz-Teil (ersten $x$-Bits = 1 setzen)
          \end{itemize}
        \item Präfix = Adresse AND Subnetzmaske
      \end{itemize}
      Subnetz
      \begin{itemize}
        \item Interfaces mit dem selben Subnetz-Teil der Adresse bilden ein Subnetz
      \end{itemize}
      Zuteilung von IP-Adressen
      \begin{itemize}
        \item Aufgabe der ICANN
          \begin{itemize}
            \item Internet Corporation for Assigned Names und Numbers
            \item Zeintraler Vertrauensanker
              \begin{itemize}
                \item Allokiert Adressen
                \item Verwaltet DNS
                \item Weist Domänennamen zu
              \end{itemize}
          \end{itemize}
        \item Delegierung der Vergabe an lokale Organisationen
          \begin{itemize}
            \item Regional IP Address Registers (RIR)
          \end{itemize}
        \item Lokale Organisationen verteilen an Provider und Unternehmen
      \end{itemize}
    \subsubsection{Das Protokoll IP}
      Internet Protocoll (IP)
      \begin{itemize}
        \item "Das" Protokoll im Internt
          \begin{itemize}
            \item Best-Effort Vermittlungsdienst
          \end{itemize}
        \item Verantwortlich für das Weiterleiten von IP-Datagrammen im Internet
      \end{itemize}
      Format eines IPv4-Datagramms\\\\
      \includegraphics[width=\textwidth]{img/ipv4_datagramm.png}
      \begin{itemize}
        \item Minimale Länge IP-Kopf: 20 Bits
      \end{itemize}
      Fragmentieren und Reassemblieren
      \begin{itemize}
        \item Anpassung an verschieden lange maximale Rahmenlängen unterliegender Netze (MTU:
          Maximum Transfer Unit)
        \item Vorgehensweise
          \begin{itemize}
            \item Flag-Felder des IP-Kopfs nutzen
              \begin{itemize}
                \item Bit 0: reserviert, muss 0 sein
                \item Bit 1:
                  \begin{itemize}
                    \item 0 = darf fragmentiert werden
                    \item 1 = darf nicht fragmentiert werden
                  \end{itemize}
                \item Bit 2:
                  \begin{itemize}
                    \item 0 = letztes Fragment
                    \item 1 = es folgen weitere Fragmente
                  \end{itemize}
                \item Fragment-Offset: Stelle, an der empfangenes Fragment in ursprüngliches
                  Datagramm eingesetzt werden muss (Basiseinheit: 8 Bytes)
              \end{itemize}
          \end{itemize}
      \end{itemize}
      Weiterleitung in IP
      \begin{itemize}
        \item Für ein Endsystem
          \begin{itemize}
            \item Rechner mit Zieladresse ist im gleichen Subnetz
              \begin{itemize}
                \item IP-Datagramm wird direkt zu diesem Empfänger geschickt
              \end{itemize}
            \item Ansonsten
              \begin{itemize}
                \item IP-Datagramm wird an voreingestelltem Router (Default-Router) weitergeleitet
              \end{itemize}
          \end{itemize}
        \item Für einen Router. Weiterleitungstabelle u.a. mit
          \begin{itemize}
            \item Angaben zum Ziel (Zieladresse oder Netzpräfix)
            \item Angabe des "Next-Hop" Routers
            \item Angabe des Interfaces, über welches das IP-Datagramm weitergeleitet werden soll
          \end{itemize}
      \end{itemize}
      Empfang eines IP-Datagramms
      \begin{itemize}
        \item Folgende Überprüfunungen werden durchgeführt
          \begin{itemize}
            \item Korrekte Länge des Kopfs?
            \item IP-Versionsnummer?
            \item Korrekte Datagrammlänge?
            \item Prüfsumme?
            \item Lebenszeit?
            \item Protokoll-Identifikation?
          \end{itemize}
        \item Falls ein Fehler erkannt wird
          \begin{itemize}
            \item Benachrichtigung mittels ICMP (Internet Control Message Protocoll)
              \begin{itemize}
                \item D.h. ICMP wird hier als Funktion aus IP ausgerufen
                \item Reagiert möglicherweise mit dem Aussenden eines ICMP-Datagramms
              \end{itemize}
          \end{itemize}
      \end{itemize}
    \subsubsection{IPv6}
      \begin{itemize}
        \item Ursprüngliche Motivation
          \begin{itemize}
            \item Größerer Adressraum durch längere Adressen: IPv6-Adressen (128 Bit)
            \item Adressraum von IPv4-Adressen (32 Bit) wird knapp bzw. geht aus $\ra$ heute kaum noch
              freie IPv4 Adressen
          \end{itemize}
        \item IPv4-Adressraum mittlerweile stark fragmentiert
          \begin{itemize}
            \item Viele kleine, verteilte Subnetze
            \item Gro0e Routingtabellen
          \end{itemize}
          $\ra$ IPv6 Adressvergabe an Netztopologie orientieren erlaubt starkes aggregieren
        \item Vereinfachungen im Vergleich zu IPv4
          \begin{itemize}
            \item Keine Fragmentierung in Routern
            \item Keine Berechnungen einer Prüfsumme
              \begin{itemize}
                \item UDP Prüfsumme jetzt verpflichtend
              \end{itemize}
            \item Einfacheres Handling der IP-Datagramme
              \begin{itemize}
                \item Einfacheres Kopfformat
                \item Feld "Header Length" nicht mehr erforderlich
                \item Kopferweiterungen statt Optionen
              \end{itemize}
            \item Einfachere Klassifikation von Paketen in Datenströme
              \begin{itemize}
                \item Einführung des Felds "Flow Label"
              \end{itemize}
          \end{itemize}
          \includegraphics[width=\textwidth]{img/ipv6_datagramm.png}
      \end{itemize}
      Typen von IPv6-Adressen
      \begin{itemize}
        \item Unicast
          \begin{itemize}
            \item Identifiziert ein dediziertes Interface
            \item Datagramm wird zu diesem Interface ausgeliefert
          \end{itemize}
        \item Multicast
          \begin{itemize}
            \item Identifiziert eine Gruppe von Interfaces
            \item Datagramm wird zu allen Mitgliedern dieser Gruppe ausgeliefert
          \end{itemize}
        \item Anycast
          \begin{itemize}
            \item Identifiziert eine Gruppe von Interfaces
            \item Datagramm wird an ein Interface der Gruppe ausgeliefert
          \end{itemize}
      \end{itemize}
      Global Unicast IPv6 Adressen
      \begin{itemize}
        \item Global routbare IPv6-Adressen
          \begin{itemize}
            \item Äquivalent zu "public IPv4 adress"
          \end{itemize}
        \item Strukturierung:
          \begin{itemize}
            \item Globales Routing Präfix
              \begin{itemize}
                \item Präfix oder "network portion"
              \end{itemize}
            \item Subnetz ID
              \begin{itemize}
                \item Bilden von Subnetzen
              \end{itemize}
            \item Interface ID
              \begin{itemize}
                \item Identifiziert Interface in Subnetz
                \item Meistens 64 Bit lang
              \end{itemize}
          \end{itemize}
      \end{itemize}
  \subsection{Kontrollebene}
    Modellierung eines Rechnernetztes als Graph
    \begin{itemize}
      \item Router (Zwischensysteme) sind Knoten
      \item Übertragungsabschnitte sind Kanten
        \begin{itemize}
          \item "Kosten" der Kanten bspw. Verzögerung, aktuelle Last, Preis ...
          \item Weg (Pfad) ist Folge von Knoten $(n_1,...,n_k)$, wobei $(n_1,n_2),(n_2,n_3),...$ Kanten
            des Graphs sind und kein Knoten mehrfach vorkommt
        \end{itemize}
    \end{itemize}
    Routingprotokoll (Finden eines "guten" Wegs
    \begin{itemize}
      \item Typischerweise ist die der Weg mit den geringsten Kosten
      \item Wie Kosten messen?
        \begin{itemize}
          \item Verschiedene Metriken denkbar
            \begin{itemize}
              \item Geringste Anzahl von Hops
              \item Schnellster Weg
              \item Am wenigsten belastbarer Weg
              \item Weg mit geringsten monetären Kosten
              \item ...
            \end{itemize}
        \end{itemize}
      \item Routingalgorithmen
    \end{itemize}
    \subsubsection{Grundlegende Routingverfahren}
      Wie dynamisch ist das Routing-Verfahren?
      \begin{itemize}
        \item Nicht adaptiv (statisch)
          \begin{itemize}
            \item Routen ändern sich nur sehr selten
            \item Routenänderungen sind viel seltener als Verkehrsänderungen
          \end{itemize}
        \item Adaptiv (dynamisch)
          \begin{itemize}
            \item Routen ändern sich in Abhänigkeit des Verkehrs bzw. der Netztopologie
              \begin{itemize}
                \item Aktueller Zustand des Netzes wird damit berücksichtigt
                \item Schleifen und Oszillation in Routen wahrscheinlicher als bei nicht adaptiven
                  Verfahren
              \end{itemize}
            \item Können periodisch operieren oder in direkter Reaktion auf Änderungen
            \item Zielkonflikt
              \begin{itemize}
                \item Router haben veraltetete oder unvollständige Informationen über den Zustand
                  des Netzes
                \item Evt. hohe Belastung durch Austausch von Routinginformationen
              \end{itemize}
          \end{itemize}
      \end{itemize}
      Routing: zentral vs. dezentral
      \begin{itemize}
        \item Zentral
          \begin{itemize}
            \item Wegewahl wird durch eine zentrale Komponente durchgeführt
              \begin{itemize}
                \item Weiterleitungstabelle in den Routern müssen von dieser zentralen Komponente
                  erstellt und verteilt werden
              \end{itemize}
            \item Zentrale Komponente benötigt globale Sicht auf Netz
          \end{itemize}
        \item Dezentral
          \begin{itemize}
            \item Routingentscheidungen/Wegewahl wird dezentral in den Routern durchgeführt
            \item Erfordert keine zentrale Komponente
            \item Isoliert
              \begin{itemize}
                \item Router entscheiden nur auf der Basis von Informationen die er selbst sammelt,
                  kein Austausch mit anderen Routern
              \end{itemize}
            \item Verteilt
              \begin{itemize}
                \item Router tauschen Routinginformationen mit Nachbarn aus
              \end{itemize}
          \end{itemize}
      \end{itemize}
      Grundlegende Routingverfahren: Überblick
      \begin{itemize}
        \item Statisches Routing
          \begin{itemize}
            \item Statische Weiterleitungstabelle
          \end{itemize}
        \item Zentralisiertes Routing
          \begin{itemize}
            \item Routing Control Center
          \end{itemize}
        \item Dezentrales Routing
          \begin{itemize}
            \item Isoliert
              \begin{itemize}
                \item Fluten
                \item Hot Potato
              \end{itemize}
            \item Verteilt
              \begin{itemize}
                \item Distanz-Vektor
                \item Link-State
              \end{itemize}
          \end{itemize}
      \end{itemize}
      Zentralisiertes Routing
      \begin{itemize}
        \item Einordnung
          \begin{itemize}
            \item Zentral, adaptiv
          \end{itemize}
        \item Zentrales Routing Control Center (RCC)
          \begin{itemize}
            \item Jeder Router sendet periodisch Zustandsinformationen an RCC
              \begin{itemize}
                \item Z.B. Liste aller aktiven Nachbarn, aktuelle Warteschlangenlänge, Umfang an
                  Verkehr seit dem letzten Bericht
              \end{itemize}
            \item RCC berechnet mit diesen Informationen optimale Wege
            \item RCC verteilt neu Weiterleitungsinformationen an Router
              \begin{itemize}
                \item Jeder Router trifft Weiterleitungsentscheidungen anhand dieser Informationen
              \end{itemize}
          \end{itemize}
      \end{itemize}
      Zentralisiertes Routing - Vor- und Nachteile
      \begin{itemize}
        \item Vorteile
          \begin{itemize}
            \item RCC hat theoretisch die vollständige Übersicht
            \item Router müssen keine aufwendige Routingberechnung durchführen
          \end{itemize}
        \item Nachteile
          \begin{itemize}
            \item Für große Netze dauert die Berechnung u.U. sehr lange
            \item Ausfall des RCC lähmt das gesamte Netz
              \begin{itemize}
                \item Erfordert daher möglichst robuste Auslegung (z.B. durch Redundanz, etwa 
                  Backup-Rechner, etc.
              \end{itemize}
            \item Inkonsistenzen möglich
              \begin{itemize}
                \item Router nahe dem RCC verfügen früher über aktualisierte Weiterleitungstabellen
                  als weiter entfernte
              \end{itemize}
            \item Starke Belastung des RCC
              \begin{itemize}
                \item Durch zentrale Funktion sowie der Links in der Nähe des RCC
              \end{itemize}
          \end{itemize}
      \end{itemize}
      Fluten
      \begin{itemize}
        \item Einordnung
          \begin{itemize}
            \item Isoliert, nicht adaptiv
          \end{itemize}
        \item Grundlegende Vorgehensweise
          \begin{itemize}
            \item Jedes eingehende Datagramm wir auf jedem Interface weitergeleitet, außer auf
              demjenigen, auf dem es empfangen wurde
          \end{itemize}
        \item Maßnahmen zur Eindämmung der Flut
          \begin{itemize}
            \item Erkennung von Duplikaten durch Sequenznummern
            \item Kontrolle der Lebensdauer eines Datagramms durch Zählen der zurückgelegten
              Übertragungsabschnitte (Hops)
              \begin{itemize}
                \item Hop-Zähler wird mit maximal erlaubter Weglänge initialisiert
                \item In jedem Router wird Zähler um 1 dekrementiert
                \item Falls Zähler den Wert 0 erreichen, muss Datagramm verworfen werden
              \end{itemize}
          \end{itemize}
        \item Varianten
          \begin{itemize}
            \item Selektives Fluten
              \begin{itemize}
                \item Weiterleitung nicht auf allen, sondern nur auf einigen Übertragungsabschnitten
              \end{itemize}
            \item Random Walk
              \begin{itemize}
                \item Zufällige Auswahl eines/mehrerer Übertragungsabschnitte(s)
              \end{itemize}
          \end{itemize}
      \end{itemize}
      Hot Potato
      \begin{itemize}
        \item Einordnung
          \begin{itemize}
            \item Isoliert, adaptiv
          \end{itemize}
        \item Jedes System versucht, eingehende Datagramme so schnell wie möglich weiterzuleiten
          \begin{itemize}
            \item Wählt Interface mit kürzester Warteschlange
            \item Weiterleitung nie auf dem Übertragungsabschnitt, auf dem Datagramm empfangen wurde
          \end{itemize}
      \end{itemize}
      Verteiltes adaptives Routing
      \begin{itemize}
        \item Systeme tauschen Routinginformationen mit Nachbarn aus
        \item Jedes System hat Routingtabelle (Informationen über zu erreichende Systeme, z.B.)
          \begin{itemize}
            \item Bevorzugert Übertragungsabschnitt zum Ziel
            \item Schätzung über Zeit oder Entfernung zum Ziel, z.B.
              \begin{itemize}
                \item Anzahl Hops
                \item Geschätzte Verzögerung in Millisekunden
                \item Geschätzte Anzahl von Datagrammen, die entlang des Weges warten
              \end{itemize}
          \end{itemize}
        \item Schätzungen werden gewonnen aus
          \begin{itemize}
            \item Zeit oder Entfernung zu den Nachbarn
              \begin{itemize}
                \item z.B. aus speziellen Echo-Paketen mit Zeitstempeln
              \end{itemize}
            \item Schätzung der Nachbarn
          \end{itemize}
        \item Varianten
          \begin{itemize}
            \item Periodischer Austausch von Routinginformationen
            \item Austausch nur bei signifikanten Änderungen
          \end{itemize}
      \end{itemize}
      Verteiltes adaptives Routing: Routingalgorithmen
      \begin{itemize}
        \item Distanz-Vektor-Algorithmen
          \begin{itemize}
            \item Routing-Metrik: Distanz
            \item Jeder Router kennt Distanz zu allen anderen Systemen im Netz
              \begin{itemize}
                \item Hierzu: Austausch der aktuellen Distanz zwischen den Nachbarn
              \end{itemize}
            \item Problem
              \begin{itemize}
                \item Kürzerer langsamerer Weg wird längerem schnellerem Weg vorgezogen
              \end{itemize}
            \item Beispielprotokolle
              \begin{itemize}
                \item Routing Information Protocol (RIP), Distance Vector Routing Protocol (DVRP)
              \end{itemize}
          \end{itemize}
        \item Link-State-Algorithmen
          \begin{itemize}
            \item Unterschiedliche Routing-Metriken pro Netz möglich
              \begin{itemize}
                \item Beispiel: Finanzielle Kosten, Auslastung
                \item Kann aktuelle Zustände der Netzinterfaces berücksichtigen
              \end{itemize}
            \item Jeder Router kennt die komplette Netztopologie
              \begin{itemize}
                \item Berechnet auf dieser Basis (konsistente) Routinginformationen
              \end{itemize}
            \item Konvergieren im Allg. schneller als Distanz-Vektor-Algorithmen
              \begin{itemize}
                \item Damit potentiell besser für größere Netzte geeignet
              \end{itemize}
            \item Beispielprotokolle
              \begin{itemize}
                \item Open Shortets Path First (OSPF)
                \item Intra-Domain Intermediate System to Intermediate System Routing (IS-IS)
              \end{itemize}
          \end{itemize}
      \end{itemize}
    \subsubsection{Intra-Domain-Routing vs. Inter-Domain-Routing}
      Internet: Netz aus Netzen
      \begin{itemize}
        \item Routing vom Endsystem A zu Endsystem B
          \begin{itemize}
            \item Von A zum lokalen Default Router
            \item Über mehrere Domänen hinweg zur Zieldomäne (Inter-Domain-Routing) (z.B. BGP)
            \item Innerhalb jeder Domäne vom Eingangs- zum Ausgangsrouter (Intra-Domain-Routing)
              (z.B. OSPF, IS-IS)
            \item In Zieldomäne zu B
          \end{itemize}
      \end{itemize}
    \subsubsection{Distanz-Vektor-Routing}
      \begin{itemize}
        \item Eigenschaften
          \begin{itemize}
            \item Verteilt
              \begin{itemize}
                \item Jeder Router enthält Informationen von seinen direkten Nachbarn, führt eine
                  Berechnung durch und verteilt dann neue Information an seine direkten Nachbarn
              \end{itemize}
            \item Iterativ
              \begin{itemize}
                \item Das Verteilen und Berechnen von Informationen geht so lange vor sich bis keine
                  neue Information mehr ausgetauscht wird
              \end{itemize}
          \end{itemize}
        \item Distanz-Vektor-Tabelle
          \begin{itemize}
            \item Grundlegende Datenstruktur für Distanz-Vektor-Algorithmen
              \begin{itemize}
              \item In jedem Router vorhanden
              \item Zeile für jedes mögliche Ziel
              \item Spalte für jeden direkten Nachbarn (Next Hop)
              \end{itemize}
            \item Weiterleiten von Daten
          \end{itemize}
      \end{itemize}
      Distanz-Vektor-Algorithmus
      \begin{itemize}
        \item Initialisierung
          \begin{itemize}
            \item $\forall\text{ Nachbarn }v: D^X(*,v)=\infty,D^X(v,v)=c(X,v)$
            \item $\forall$ Ziele $y:$ sende $\min_wD^X(y,w)$ zu jedem Nachbarn, wobei $w$ alle
              Nachbarn enthält
          \end{itemize}
        \item Schleife
          \begin{itemize}
            \item Geänderte Übertragungsabschnittskosten
              \begin{itemize}
                \item $c(X,V$ ändert sich um Wert $d$ (kann positiv oder negativ sein)
                \item $\forall$ Ziele $y$: $D^X(y,V)=D^X(y,V)+d$
              \end{itemize}
            \item Update-Nachricht von einem Nachbarn
              \begin{itemize}
                \item Kürzester Pfad von $V$ zu einem Ziel $Y$ hat sich geändert zu "neuer Wert"
                \item $D^X(Y,V)=c(X,V)+$ "neuer Wert" für dieses Zeil
              \end{itemize}
            \item Falls neuer $\min_wD^X(Y,w)$ für ein Ziel $Y$ existiert, sende diesen Wert zu allen
              Nachbarn
          \end{itemize}
        \item Betrachteter Algorithmus: Bellman-Ford-Algorithmus
          \begin{itemize}
            \item Eingesetzt im Routingprotokoll RIP (Routing Information Protocol)
          \end{itemize}
      \end{itemize}
    \subsubsection{Link-State-Routing}
      \begin{itemize}
        \item Grundlegende Vorgehensweise
          \begin{itemize}
            \item Systeme kennen am Anfang nur ihre direkten Nachbarn
            \item Entdecken neuer Nachbarn mittels spezieller Routingnachrichten
              \begin{itemize}
                \item z.B. HELLO-Nachricht
              \end{itemize}
            \item Bestimmen der Kosten zu direkten Nachbarn
            \item Link State Broadcast
              \begin{itemize}
                \item Identität und Kosten zu direkten Nachbarn werden an alle Router im Netz
                  weitergeleitet (Fluten)
                \item Systeme können durch empfangene Link State Broadcasts Topologie lernen
                \item Ergebnis: Alle Systeme haben identisches Wissen über das Netz
              \end{itemize}
          \end{itemize}
        \item Berechnung der kürzesten Pfade durch Link-State-Algorithmus
          \begin{itemize}
            \item Jedes System berechnet de kürzesten Pfade
            \item Die berechneten Pfade sind aufgrund der identischen Information gleich
          \end{itemize}
        \item Nach Fluten und Berechnung der kürzesten Pfade in jedem System
          \begin{itemize}
            \item Schleifenfrei
            \item In stabilen Zustand konvergiert
          \end{itemize}
        \item Im folgenden betrachteter Algorithmus: Dijkstra-Algorithmus
          \begin{itemize}
            \item Berechnet Pfade mit geringsten Kosten von einem System zu allen anderen Systemen
              im Netz
          \end{itemize}
        \item Eingesetzt in den Routingprotokollen
          \begin{itemize}
            \item OSPF (Open Shortest Path First) 
            \item IS-IS (Intermediate System to Intermediate System)
          \end{itemize}
      \end{itemize}
      \includegraphics[width=\textwidth]{img/link_state.png}
      Link-State vs. Distanz-Vektor
      \begin{itemize}
        \item Link-State
          \begin{itemize}
            \item Konvergiert schnell
              \begin{itemize}
                \item Link-State-Nachrichten können während der Berechnung verschickt werden
                \item Fluten ist unabhängig von berechneten Pfaden, unabhängig von Routing-Schleifen
              \end{itemize}
            \item Kompliziert zu implementieren
              \begin{itemize}
                \item Speichern aller empfangenen Link-State-Nachrichten
                \item Rekonstruktion der Topologie erforderlich
              \end{itemize}
          \end{itemize}
        \item Distanz-Vektor
          \begin{itemize}
            \item Konvergiert langsam
              \begin{itemize}
                \item Bei "schlechten Nachrichten": Schleifen und Count-to-Infinity-Probleme
                \item Bei "guten Nachrichten": Routen müssen nacheinander pro Router angepasst werden
              \end{itemize}
            \item Leicht zu implementieren
              \begin{itemize}
                \item Update der Tabelleen beim Empfang einer Nachricht
              \end{itemize}
          \end{itemize}
        \item Gewinner?
          \begin{itemize}
            \item Distanz-Vektor: Gut für kleine Netze
            \item Link-State: Besser geeignet für größere (aber nicht sehr große) Netze
          \end{itemize}
      \end{itemize}
  \subsection{Weitere Protokolle der Kontrollebene}
    \subsubsection{Zuteilung von Adressen: DHCP}
      Zuteilung von IP-Adressen
      \begin{itemize}
        \item Manuelle Konfiguration
          \begin{itemize}
            \item Durch Systemadministrator
              \begin{itemize}
                \item Windows: Systemsteuerung $\ra$ Netzwerkstatus und Aufgaben anzeigen $\ra$
                  Adaptereinstellungen
                \item Unix: /etc/rc.config
                \item Linux: /etc/network/interfaces, /etc/sysconfig, ...
              \end{itemize}
          \end{itemize}
        \item Dynamische Konfiguration
          \begin{itemize}
            \item Dynamische Host Configuration Protocol (DHCP)
            \item DHCP-Server liefert bei Anfrage eine IP-Adresse an den Cleient
              \begin{itemize}
                \item "plug-and-play", "zero-config"
              \end{itemize}
          \end{itemize}
      \end{itemize}
      DHCP: Dynamic Host Configuration Protocol
      \begin{itemize}
        \item Aufgabe
          \begin{itemize}
            \item Dynamischer Bezug von IP-Adressen durch Endsysteme
            \item Anforderung einer IP-Adresse bei Beitritt zum Subnetz
          \end{itemize}
        \item Eigenschaften
          \begin{itemize}
            \item Client-Server-Protokoll
              \begin{itemize}
                \item Zuteilung von IP-Adressen durch DHCP-Server
              \end{itemize}
            \item Zeitliche Gültigkeit der Zuteilung beschränkt
              \begin{itemize}
                \item Kann auf Anfrage verlängert werden
                \item Wiedervergabe nach Ablauf möglich
              \end{itemize}
            \item Unterstützt mobile Endsysteme
              \begin{itemize}
                \item Aber keine "seamless mobility": neue IP-Adresse in neuem Subnetz
                \item Bei Wechsel der IP-Adresse $\ra$ Unterbrechung einer aktiven TCP-Verbindung
              \end{itemize}
          \end{itemize}
        \item Funktionsweise für IPv4
        \item Standardablauf
          \begin{itemize}
            \item Endsystem tritt einem Subnetz bei
              \begin{itemize}
                \item Endsystem sind keine IP-Adressen von DHCP-Servern bekannt
              \end{itemize}
            \item Austausch von DHCP-Nachrichten zwischen Client und Server
              \begin{itemize}
                \item Endsystem sendet Brodcast: "DHCP discover"
                \item DHCP-Server antwortet mit: "DHCP offer"
                \item Endsystem fordert IP-Adresse an: "DHCP request"
                \item DHCP-Server teilt IP-Adresse zu: "DHCP ack"
              \end{itemize}
          \end{itemize}
      \end{itemize}
      DHCP: Mehr als nur IPv4-Adressen
      \begin{itemize}
        \item Mit DHCP können weitere Informationen im Subnetz verteilt werden
          \begin{itemize}
            \item Adressen des Default-Routers für das Subnetz
            \item Name und IP-Adresse des DNS-Servers
            \item Subnetzmaske
            \item Zeitserver (NTP - Network Time Protocol)
          \end{itemize}
      \end{itemize}
    \subsubsection{ICMP}
      Internet Control Message Protocol (ICMP)
      \begin{itemize}
        \item Einzelne Verluste von Datagrammen werden von IP nicht gemeldet
          \begin{itemize}
            \item IP stellt zuverlässigen Dienst bereit
          \end{itemize}
        \item Schwerwiegende Probleme (z.B. Unterbrechung einer Leitung)
          \begin{itemize}
            \item werden Kommunikationspartnern zur Vermeidung von Folgefehlern mittels ICMP mitgeteilt
          \end{itemize}
        \item ICMP unterstützt den Austausch von Fehlernachrichten, Statusanfrage und
          Zustandsinformationen
      \end{itemize}
      ICMP-Statusanfrage
      \begin{itemize}
        \item Echo und Echoantwort (echo and echo reply)
          \begin{itemize}
            \item Dient der Überprüfung der Aktivität von Geräten
            \item Der Empfänger einer Echo-Anfrage sendet in der Echo-Antwort die erhaltenen Daten an
              den Kommunikationspartner zurück
          \end{itemize}
        \item Zeitstempel und Zeitstempelantwort (timestamp and timestamp reply)
          \begin{itemize}
            \item Dient der Bestimmung von Umlaufzeiten (Round Trip Time, RTT)
            \item Datagramme umfassen mehrere Felder zur Aufnahme von Zeitstempeln
              \begin{itemize}
                \item Damit: Berechnung der Bearbeitungszeiten beim Empfänger und der Verzögerung im
                  Netz
              \end{itemize}
          \end{itemize}
      \end{itemize}
      ICMP-Datagramme
      \begin{itemize}
        \item Übertragung
          \begin{itemize}
            \item Im Datenteil von IP-Datagrammen
            \item Protocol-Feld im Kopf hat Wert "1"
          \end{itemize}
        \item Format ICMP-Datagramm
          \begin{itemize}
            \item Type: Typ des Datagramms (z.B. Type = 3 entspricht "Zieladresse nicht erreichbar"
            \item Code: Genaue Beschreibung des Datagramms (z.B. "Netz nicht erreichbar")
            \item Checksum: Prüfsumme über gesamtes ICMP-Datagramm
            \item Info: abhängig von Typ des ICMP-Datagramms (z.B. Felder für Zeitstempel bei
              "Zeitstempel und Zeitstempelantwort"
          \end{itemize}
      \end{itemize}
    \subsubsection{Vermittlungsschicht im Internet}
      \begin{itemize}
        \item Sowol im Endsystem als auch im Router vorhanden
      \end{itemize}
      Aufgaben der Protokolle
      \begin{itemize}
        \item TCP (Transmission Control Protocol)
          \begin{itemize}
            \item Stellt zuverlässigen Transportdienst bereit
          \end{itemize}
        \item UDP (User Datagramm Protocol)
          \begin{itemize}
            \item Stellt unzuverlässigen Transportdienst bereit
          \end{itemize}
        \item IP (Internet Protocol)
          \begin{itemize}
            \item Unzuverlässige Übertragung/Weiterleitung von IP-Datagrammen
          \end{itemize}
        \item ICMP (Internet Control Message Protocol)
          \begin{itemize}
            \item Austausch von Kontrollinformationenen innerhalb der Vermittlungsschicht
          \end{itemize}
        \item ARP (Address Resolution Protocol)
          \begin{itemize}
            \item Zuordnung von IP-Adressen zu Adressen der Sicherungsschicht (MAC-Adressen)
          \end{itemize}
        \item RARP (Reverse Address Resolution Protocol)
          \begin{itemize}
            \item Stellt Umkehrfunktion von ARP zur Verfügung
          \end{itemize}
        \item Routingprotokolle
          \begin{itemize}
            \item z.B. RIP (Routing Information Protocol), OSPF (Open Shortest Path First), BGP
              (Border Gateway Protocol)
          \end{itemize}
      \end{itemize}
    \subsection{Software-Defined Networking}
      Software-Defined Networking (SDN)
      \begin{itemize}
        \item Über vier grundlegende Eigenschaften (E1-E4)
      \end{itemize}
      \includegraphics[width=\textwidth]{img/sdn.png}\\
      Traditionelles IP-Routing
      \begin{itemize}
        \item Kontroll- und Datenebene in jedem Router
        \item Vorteile
          \begin{itemize}
            \item Ausfallsicherheit
              \begin{itemize}
                \item Selbstorganisiert
                \item Verteilte Kontrolle, hohe Redundanz
              \end{itemize}
            \item Schnelle Reaktion
              \begin{itemize}
                \item Hochoptimierte Routinghardware
                \item Lokale Routingentscheidung
              \end{itemize}
            \item Bewährtes Konzept
          \end{itemize}
        \item Nachteile
          \begin{itemize}
            \item Proprietär
              \begin{itemize}
                \item Herstellerspezifische Management-Schnittstellen
                \item Mischbetrieb schwierig
              \end{itemize}
            \item Unflexibel
              \begin{itemize}
                \item Einführung neuer Funktion schwierig/langwierig
                \item Aufwändige Standardisierungsprozesse
              \end{itemize}
            \item Teuer
              \begin{itemize}
                \item Hochqualifiziertes Personal benötigt
                \item Overprovisioning (für Lastspitzen)
              \end{itemize}
          \end{itemize}
      \end{itemize}
      \pagebreak
      SDN:Kontrollebene in logisch zentralem Kontroller\\\\
      \includegraphics[width=\textwidth]{img/sdn_controller.png}\\\\
      Traditionell vs. SDN
      \begin{itemize}
        \item Traditioneller IP-Router
          \begin{itemize}
            \item Kennt keine Flows
            \item Weiterleitung über Ziel-IP-Adresse $Longest Prefix Matching$
          \end{itemize}
        \item SDN-fähiger Switch
          \begin{itemize}
            \item Weiterleitung über Flowtabelle
            \item Nutzt verschiedene Felder im IP-Kopf
            \item Speichert Zustand pro Flow
          \end{itemize}
      \end{itemize}
      Proaktive Flow-Programmierung
      \begin{itemize}
        \item Controller stellt sicher, dass immer ein passender Eintrag vorhanden ist
          \begin{enumerate}[label=\arabic*.]
            \item Controller: programmiert passenden Flow auf Switch
            \item Endsystem 1: sendet Paket
            \item Switch: Packet-in (Port 1)
            \item Switch: Flow-Table Lookup
            \item Macht $\ra$ \verb!Action: "Output:2"!
            \item Switch: Paketweiterleitung $\ra$ (von Port 1 nach Port 2)
          \end{enumerate}
      \end{itemize}
      Reaktive Flow-Programmierung
      \begin{itemize}
        \item Fall 1: Kein passender Eintrag in Flow-Table vorhanden
          \begin{enumerate}[label=\arabic*.]
            \item Endsystem 1: Sendet Paket
            \item Switch: Packet-in (Port 1)
            \item Switch: Flow-Table Lookup
            \item Mismatch
            \item Switch: Weiterleitung zum Controller
            \item Controller: Programmiert passenden Flow auf Switch
            \item Controller: Explizite \verb!Packet-Out! Nachricht an Switch
            \item Switch: Sendet Paket über Port 2
          \end{enumerate}
        \item Fall 2: Passender Eintrag vorhanden
          \begin{enumerate}[label=\arabic*]
            \item Endsystem 1: Sendet Paket
            \item Switch: Packet-in (Port 1)
            \item Switch: Flow-Table Lookup
            \item Match $\ra$ \verb!Action: "Output:2"!
            \item Switch: Paketweiterleitung $\ra$ (von Port 1 nach Port 2)
          \end{enumerate}
      \end{itemize}
      \includegraphics[width=\textwidth]{img/reaktie_flow_programmierung.png}\\\\
      SDN: Pro und Contra
      \begin{itemize}
        \item Pro
          \begin{itemize}
            \item Logisch zentralisierte Sicht
              \begin{itemize}
                \item Controller hat Überblick über das gesamte Netz
                \item Einfacher Einsatz von Graphen-Algorithmen
              \end{itemize}
            \item Neue Funktionalität in Software
              \begin{itemize}
                \item Als App im Controller
                \item Deutlich kürzere Entwicklungszyklen
              \end{itemize}
            \item Trennung von Kontrollebene und Datenenbene
              \begin{itemize}
                \item Innovation unabhängig voneinander möglich
                \item Herstellerunabhängigkeit durch offene Schnittstelle
              \end{itemize}
          \end{itemize}
        \item Contra
          \begin{itemize}
            \item Skalierbarkeitn
            \item "Single Point of Failure"
            \item Kommunikation mit Controller langsam im Vergleich zum Datenpfad
            \item Anzahl TCAM Einträge begrenzt
          \end{itemize}
      \end{itemize}
